% !TeX root = ../main.tex

\xchapter{燃料力学非线性、裂变气体与间隙传热模型的耦合实现及验证}{Coupled Implementation and Verification of Fuel Mechanical Nonlinearity, Fission Gas, and Gap Heat Transfer Models}

\xsection{引言}{Introduction}
第二章建立了氧化物燃料元件热学、力学、微观结构与组分演化的多物理场控制方程及总体耦合框架，第三章进一步给出了MOX燃料与包壳材料行为、气孔迁移、Pu/Am再分布及O/M演化等闭合模型。在此基础上，还需要将不同物理模型转化为适用于燃料性能分析程序的数值算法，并解决非线性本构积分、裂变气体状态演化以及燃料--包壳间隙传热之间的数据传递与耦合问题。

燃料元件热--力耦合计算中，塑性、蠕变等非弹性材料行为需要在材料积分点进行本构积分，更新应力及内部变量，并与整体准静态平衡方程耦合求解；热膨胀和辐照肿胀等则按相应本征应变模型处理。裂变气体行为则涉及晶内扩散、晶界聚集、气泡演化及释放等过程，其计算结果进一步影响棒内气体组成、气腔压力及燃料--包壳间隙传热状态。因此，局部本构积分、裂变气体演化与间隙传热模型之间的稳定耦合，是实现燃料元件多物理场数值分析的关键环节。

本章重点研究上述模型的数值实现与程序耦合方法。首先，基于$J_2$塑性理论和隐式径向返回算法，实现燃料与包壳非线性力学本构的积分更新，并介绍其在MOOSE材料计算框架中的组织方式；随后，分别实现Forsberg--Massih裂变气体模型和SCIANTIX细观机理模型，建立裂变气体释放、棒内气体状态与间隙传热之间的数据传递关系；最后，通过单元级、子模型级及耦合算例对相关算法和程序实现进行验证，为后续燃料棒尺度多物理场分析提供数值基础。

本章所采用的J2塑性理论、径向返回算法及Forsberg–Massih裂变气体模型均具有已有理论基础，SCIANTIX则为公开的晶粒尺度机理程序。本文的工作重点并非重新提出上述基础模型，而是在MOOSE框架内实现非线性本构的隐式积分、裂变气体机理程序的状态变量传递及间隙传热的耦合更新，并通过分层数值测试评价相关实现的正确性与适用性。

\xsection{热弹塑性力学模型的离散与隐式积分}{Discretization and Implicit Integration of Thermo-Elasto-Plastic Mechanics}
\xsubsection{非线性力学问题的数值离散}{Numerical Discretization of the Nonlinear Mechanical Problem}
第二章已给出燃料与包壳的准静态力学平衡方程、应变分解关系及压力边界条件，第三章进一步给出了相应材料行为闭合模型。有限元弱形式、空间离散以及Newton迭代等属于非线性有限元与计算塑性的通用理论，其系统推导可参见经典计算塑性专著\cite{dunneIntroductionComputationalPlasticity2005}。因此，本节不再重复一般有限元理论，而重点给出与后续本构积分和程序实现直接相关的离散关系。

在小变形条件下，总应变可统一分解为：
\begin{equation}
	\boldsymbol{\varepsilon}
	=
	\boldsymbol{\varepsilon}^{\mathrm{e}}
	+
	\boldsymbol{\varepsilon}^{\mathrm{th}}
	+
	\boldsymbol{\varepsilon}^{\mathrm{in}}
	\label{eq:ch4_strain_decomposition}
\end{equation}
\noindent 式中：$\boldsymbol{\varepsilon}^{\mathrm{e}}$和$\boldsymbol{\varepsilon}^{\mathrm{th}}$分别为弹性应变和热应变，$\boldsymbol{\varepsilon}^{\mathrm{in}}$用于统一表示塑性应变、蠕变应变以及按本征应变处理的辐照肿胀等非热应变分量，各应变分量均为无量纲量。应力由材料本构关系确定：
\begin{equation}
	\boldsymbol{\sigma}
	=
	\mathbf{D}:
	\left(
	\boldsymbol{\varepsilon}
	-\boldsymbol{\varepsilon}^{\mathrm{th}}
	-\boldsymbol{\varepsilon}^{\mathrm{in}}
	\right)
	\label{eq:ch4_constitutive_relation}
\end{equation}
\noindent 式中：$\mathbf{D}$为弹性刚度张量（$\mathrm{Pa}$），其具体材料参数由第三章给出的材料模型确定。

对第二章中的准静态力学平衡方程进行有限元离散后，第$i$个自由度对应的残差可写为：
\begin{equation}
	R_i(\mathbf{u})
	=
	\int_{\Omega}
	\boldsymbol{\sigma}:\nabla^{s}\boldsymbol{\varPhi}_i
	\,\mathrm{d}\Omega
	-
	\int_{\Omega}
	\mathbf{f}\cdot\boldsymbol{\varPhi}_i
	\,\mathrm{d}\Omega
	-
	\int_{\Gamma_t}
	\bar{\mathbf{t}}\cdot\boldsymbol{\varPhi}_i
	\,\mathrm{d}\Gamma
	=0
	\label{eq:ch4_mechanical_residual}
\end{equation}
\noindent 式中：$\boldsymbol{\varPhi}_i$为位移试函数（无量纲），$\nabla^{s}$表示对称梯度算子，$\mathbf{f}$为单位体积体力（$\mathrm{N\,m^{-3}}$），$\bar{\mathbf{t}}$为给定表面力（$\mathrm{N\,m^{-2}}$）。燃料棒内、外压力以及接触载荷均通过相应边界条件参与整体力学平衡求解。

全局非线性方程采用Newton迭代求解。在每次全局迭代中，由当前位移增量获得材料积分点处的应变增量，并调用相应本构积分算法更新应力及内部状态变量；根据更新后的积分点应力计算单元内力并组装整体平衡残差，同时由材料本构切线构造整体切线刚度矩阵。由此，燃料元件非线性力学求解可分为全局平衡迭代与局部本构积分两个层次。对于本文涉及的弹塑性及黏塑性材料，局部状态更新采用基于$J_2$理论的隐式径向返回算法。

\xsubsection{径向返回算法}{Radial Return Algorithm}

径向返回算法是$J_2$塑性及相应$J_2$型黏塑性本构积分中常用的隐式应力更新方法，其完整理论推导、关联流动法则及一致切线构造可参见计算塑性专著\cite{dunneIntroductionComputationalPlasticity2005}。本文不再展开von Mises屈服面的几何性质及一般Newton方法推导，仅给出与程序实现直接相关的状态更新关系。

对于采用von Mises屈服准则及各向同性硬化规律的$J_2$塑性模型，屈服函数写为：
\begin{equation}
	f
	=
	\sigma_{\mathrm{eq}}
	-
	\sigma_y\left(\bar{\varepsilon}^{\mathrm{p}}\right)
	\label{eq:ch4_j2_yield}
\end{equation}

其中，等效应力为：
\begin{equation}
	\sigma_{\mathrm{eq}}
	=
	\sqrt{\frac{3}{2}
		\boldsymbol{s}:\boldsymbol{s}}
	\qquad
	\boldsymbol{s}
	=
	\boldsymbol{\sigma}
	-\frac{1}{3}\operatorname{tr}
	(\boldsymbol{\sigma})\mathbf{I}
	\label{eq:ch4_equivalent_stress}
\end{equation}
\noindent 式中：$\boldsymbol{s}$为偏应力张量（$\mathrm{Pa}$），$\bar{\varepsilon}^{\mathrm{p}}$为累积等效塑性应变（无量纲），$\sigma_y\left(\bar{\varepsilon}^{\mathrm{p}}\right)$为当前屈服应力（$\mathrm{Pa}$）。

在时间步$k\rightarrow k+1$内，首先假定应变增量全部为弹性增量，并根据当前材料状态构造弹性试算应力$\boldsymbol{\sigma}^{\mathrm{trial}}$及其等效应力$\sigma_{\mathrm{eq}}^{\mathrm{trial}}$。相应的试屈服函数为：
\begin{equation}
	f^{\mathrm{trial}}
	=
	\sigma_{\mathrm{eq}}^{\mathrm{trial}}
	-
	\sigma_y\left(\bar{\varepsilon}^{\mathrm p,k}\right)
	\label{eq:ch4_trial_yield}
\end{equation}

当$f^{\mathrm{trial}}\leq 0$时，弹性试算应力位于屈服面内或屈服面上，当前时间步采用弹性更新：
\begin{equation}
	\boldsymbol{\sigma}^{k+1}
	=
	\boldsymbol{\sigma}^{\mathrm{trial}}
	\qquad
	\bar{\varepsilon}^{\mathrm p,k+1}
	=
	\bar{\varepsilon}^{\mathrm p,k}
	\label{eq:ch4_elastic_update}
\end{equation}

当$f^{\mathrm{trial}}>0$时，需要进行塑性修正。对于关联流动的$J_2$模型，塑性修正方向与试偏应力方向一致，塑性应变增量可写为：
\begin{equation}
	\Delta\boldsymbol{\varepsilon}^{\mathrm{p}}
	=
	\frac{3}{2}
	\Delta\bar{\varepsilon}^{\mathrm{p}}
	\frac{\boldsymbol{s}^{\mathrm{trial}}}
	{\sigma_{\mathrm{eq}}^{\mathrm{trial}}}
	\label{eq:ch4_plastic_strain_increment}
\end{equation}

其中，$\Delta\bar{\varepsilon}^{\mathrm{p}}$为累积等效塑性应变增量。采用径向返回关系后，局部本构积分可归结为关于$\Delta\bar{\varepsilon}^{\mathrm{p}}$的标量非线性方程：
\begin{equation}
	g\left(\Delta\bar{\varepsilon}^{\mathrm{p}}\right)
	=
	\sigma_{\mathrm{eq}}^{\mathrm{trial}}
	-
	3G\Delta\bar{\varepsilon}^{\mathrm{p}}
	-
	\sigma_y
	\left(
	\bar{\varepsilon}^{\mathrm p,k}
	+\Delta\bar{\varepsilon}^{\mathrm{p}}
	\right)
	=0
	\label{eq:ch4_return_mapping_residual}
\end{equation}
\noindent 式中：$G$为剪切模量（$\mathrm{Pa}$）。本文采用Newton迭代求解式\eqref{eq:ch4_return_mapping_residual}：
\begin{equation}
	\Delta\bar{\varepsilon}^{\mathrm p,(j+1)}
	=
	\Delta\bar{\varepsilon}^{\mathrm p,(j)}
	-
	\frac{
		g\!\left(
		\Delta\bar{\varepsilon}^{\mathrm p,(j)}
		\right)
	}{
		g'\!\left(
		\Delta\bar{\varepsilon}^{\mathrm p,(j)}
		\right)
	}
	\label{eq:ch4_return_newton}
\end{equation}

对于各向同性硬化模型，令当前硬化模量$h=\mathrm{d}\sigma_y/\mathrm{d}\bar{\varepsilon}^{\mathrm p}$，则：
\begin{equation}
	g'\left(\Delta\bar{\varepsilon}^{\mathrm{p}}\right)
	=
	-3G-h
	\label{eq:ch4_return_derivative}
\end{equation}

当局部残差满足给定收敛准则后，采用收敛的$\Delta\bar{\varepsilon}^{\mathrm{p}}$更新累积等效塑性应变：
\begin{equation}
	\bar{\varepsilon}^{\mathrm p,k+1}
	=
	\bar{\varepsilon}^{\mathrm p,k}
	+
	\Delta\bar{\varepsilon}^{\mathrm{p}}
	\label{eq:ch4_equivalent_plastic_strain_update}
\end{equation}

同时更新塑性应变张量：
\begin{equation}
	\boldsymbol{\varepsilon}^{\mathrm p,k+1}
	=
	\boldsymbol{\varepsilon}^{\mathrm p,k}
	+
	\Delta\boldsymbol{\varepsilon}^{\mathrm{p}}
	\label{eq:ch4_plastic_strain_update}
\end{equation}

以及应力状态：
\begin{equation}
	\boldsymbol{\sigma}^{k+1}
	=
	\boldsymbol{\sigma}^{\mathrm{trial}}
	-
	2G\Delta\boldsymbol{\varepsilon}^{\mathrm{p}}
	\label{eq:ch4_stress_update}
\end{equation}

径向返回算法的计算流程如图\ref{fig:stress_strain_update}所示。该算法在材料积分点完成试应力构造、屈服判断、局部非线性迭代和状态变量更新，并将更新后的应力及算法一致切线模量提供给整体力学求解器，从而形成“全局平衡迭代--局部本构积分”的嵌套求解过程。

\begin{figure}[!htb]
	\centering
	\includegraphics[width=13.5cm,height=9.00cm,keepaspectratio]{stress_strain_update.pdf}
	\caption{径向返回算法计算过程}
	\label{fig:stress_strain_update}
\end{figure}

\xsubsection{代码实现和程序框架}{Code Implementation and Computational Framework}

基于上述本构积分方法，本文依托MOOSE材料计算框架实现非线性应力更新。MOOSE将全局平衡方程求解与材料积分点处的局部本构更新相分离：具体的非弹性应力更新模型继承自\texttt{StressUpdateBase}，不直接参与全局残差计算，而由上层材料\texttt{Compute\allowbreak Multiple\allowbreak Inelastic\allowbreak Stress}统一调用。该组织方式使同一时间步内的局部本构模型能够根据收敛需要被重复迭代，从而实现应力与内部状态变量的稳定更新。

在每一时间步和积分点上，\texttt{Compute\allowbreak Multiple\allowbreak Inelastic\allowbreak Stress}通过\texttt{computeQpStress()}启动应力更新流程，其父类\texttt{Compute\allowbreak Multiple\allowbreak Inelastic\allowbreak StressBase}负责组织一个或多个非弹性模型的迭代调用。计算首先根据当前总应变增量构造试应力，随后由\texttt{computeAdmissibleState()}调用具体的应力更新模型完成局部本构积分。对于径向返回算法，其核心实现位于\texttt{RadialReturnStressUpdate}类中；该类继承自\texttt{StressUpdateBase}，通过标量非线性方程求解等效非弹性应变增量，并据此校正应力、更新应变分量及内部变量。

上述类之间形成由“全局材料驱动--局部状态更新--一致切线返回”构成的层次化计算结构。主要接口及其功能分别列于表\ref{tab:multi_inelastic_base}、表\ref{tab:stress_update_base}和表\ref{tab:radial_return}中，整体类关系如图\ref{fig:type_relationship}所示。

{

\begin{xltabular}{\textwidth}{@{}
	>{\centering\arraybackslash}m{4.0cm}
	>{\centering\arraybackslash}m{2.6cm}
	>{\centering\arraybackslash}X
	@{}}

	\caption{\texttt{Compute\allowbreak Multiple\allowbreak Inelastic\allowbreak Stress}的关键函数及功能}
	\label{tab:multi_inelastic_base}
	\\

	\toprule
	函数名称 & 类型 & 功能说明 \\
	\midrule
	\endfirsthead

	\multicolumn{3}{c}{
		{表~\thetable（续）\quad
				\texttt{Compute\allowbreak Multiple\allowbreak Inelastic\allowbreak Stress}的关键函数及功能}
	}
	\\
	\toprule
	函数名称 & 类型 & 功能说明 \\
	\midrule
	\endhead

	\midrule
	\multicolumn{3}{r}{续下页}
	\\
	\endfoot

	\bottomrule
	\endlastfoot

	\texttt{computeQpStress()}
	& 主流程
	& 在当前积分点组织应力计算，是整体应力更新的入口
	\\
	\addlinespace

	\texttt{updateQpState()}
	& 多模型迭代
	& 依次调用用户指定的非弹性模型，迭代更新应力和非弹性应变，并计算本构切线
	\\
	\addlinespace

	\texttt{computeAdmissibleState()}
	& 本构调用
	& 将弹性试算应力和应变增量传递给具体\texttt{StressUpdateBase}模型，完成局部状态更新
	\\
	\addlinespace

	\texttt{computeQpJacobianMult()}
	& 切线组装
	& 根据弹性张量及各非弹性模型返回的一致切线构造整体Newton迭代所需的雅可比矩阵贡献
	\\

\end{xltabular}
}

{
\begin{xltabular}{\textwidth}{@{}
	>{\centering\arraybackslash}m{4.0cm}
	>{\centering\arraybackslash}m{2.6cm}
	>{\centering\arraybackslash}X
	@{}}

	\caption{\texttt{StressUpdateBase}的关键接口及功能}
	\label{tab:stress_update_base}
	\\

	\toprule
	函数名称 & 类型 & 功能说明 \\
	\midrule
	\endfirsthead

	\multicolumn{3}{c}{
		表~\thetable（续）\quad
		\texttt{StressUpdateBase}的关键接口及功能
	}
	\\
	\toprule
	函数名称 & 类型 & 功能说明 \\
	\midrule
	\endhead

	\midrule
	\multicolumn{3}{r}{续下页}
	\\
	\endfoot

	\bottomrule
	\endlastfoot

	\texttt{updateState()}
	& 核心接口
	& 根据应变增量和试应力调用具体本构模型，返回更新后的应力、弹性应变增量和非弹性应变增量
	\\
	\addlinespace

	\texttt{updateStateSubstep()}
	& 子步积分
	& 在单个子步内执行状态更新，用于提高大应变增量或强非线性条件下的积分稳定性
	\\
	\addlinespace

	\texttt{propagate\allowbreak Qp\allowbreak Stateful\allowbreak Properties()}
	& 状态管理
	& 在未执行本构更新或时间推进过程中传递积分点历史状态变量
	\\

\end{xltabular}
}
从算法流程上看，MOOSE框架下的径向返回算法的应力更新主要包括三个阶段。首先为弹性预测阶段，即在假设材料响应完全弹性的前提下计算试应力；其次为塑性修正阶段，当试应力违反屈服条件时，通过 returnMappingSolve() 对等效非弹性应变增量进行迭代求解，实现应力状态向屈服面的径向返回；最后为状态更新阶段，在 compute\allowbreak Stress\allowbreak Finalize() 中完成应力、弹性与非弹性应变增量及内部变量的更新，并计算算法一致切线模量，为整体Newton迭代提供线性化信息，以改善非线性求解的收敛性能。表\ref{tab:radial_return}给出RadialReturnStressUpdate类中的主要函数和功能说明。

\begin{table}[H]
	\centering
	\caption{\texttt{RadialReturnStressUpdate}的关键函数及功能}
	\label{tab:radial_return}
	\begin{tabularx}{\textwidth}{
		>{\centering\arraybackslash}m{3.5cm}
		>{\centering\arraybackslash}m{2.0cm}
		>{\centering\arraybackslash}X
		}
		\toprule
		函数名称 & 类型                              & 功能说明 \\
		\midrule
		\texttt{updateState()}
		     & 核心算法
		     & 组织试应力计算、屈服判断、返回映射及非弹性状态更新              \\

		\texttt{compute\allowbreak Stress\allowbreak Initialize()}
		     & 初始化
		     & 在返回映射迭代开始前计算弹性试算应力、试算等效应力及相关内部状态           \\

		\texttt{returnMappingSolve()}
		     & 局部求解
		     & 通过单变量非线性迭代求解等效非弹性应变增量                  \\

		\texttt{compute\allowbreak Stress\allowbreak Finalize()}
		     & 状态更新
		     & 在局部迭代收敛后完成应力及内部状态变量的最终更新               \\

		\texttt{compute\allowbreak Tangent\allowbreak Operator()}
		     & 切线计算
		     & 根据收敛后的材料状态构造算法一致切线模量，用于整体Newton迭代        \\

		\texttt{calculate\allowbreak Number\allowbreak Substeps()}
		     & 子步控制
		     & 根据应变增量和给定准则确定局部积分所需的子步数量               \\
		\bottomrule
	\end{tabularx}
\end{table}

此外，为提高算法在强非线性条件下的稳定性，RadialReturnStressUpdate 还引入了子步积分（substepping）策略。当单步应变增量过大时，可将其划分为若干子步，并在每个子步内重复执行应力更新过程，从而有效控制积分误差并增强收敛性能。与此同时，Compute\allowbreak Multiple\allowbreak Inelastic\allowbreak Stress 支持多个非弹性模型的耦合计算，通过 Picard 迭代在不同本构模型之间循环更新应力与应变分量，直至满足收敛判据。这种模块化与层次化的设计，使得不同类型的非弹性行为（如塑性、蠕变及其耦合）能够在统一框架下灵活组合，并显著提升了复杂材料本构模拟的扩展性与鲁棒性。

\begin{figure}[!htb]
	\centering
	\includegraphics[width=13.5cm,height=9.00cm,keepaspectratio]{radial_return_moose.pdf}
	\caption{MOOSE框架下径向返回算法类之间的关系}
	\label{fig:type_relationship}
\end{figure}

\xsection{裂变气体模型与间隙传热模型}{Fission Gas and Gap Heat Transfer Models}
第二章已给出裂变气体、棒内气体状态与燃料--包壳间隙传热之间的总体耦合关系。本节在此基础上重点说明裂变气体模型及间隙传热模型的数值实现与程序接口，不再重复一般性的物理过程描述。

针对裂变气体行为，本文采用两种不同建模层次的方法：一是Forsberg--Massih简化机理模型，用于描述晶内扩散、晶界累积及释放等主要过程；二是耦合开源裂变气体行为程序SCIANTIX，以获得晶内气泡、晶界气体及高燃耗结构等更细致的状态信息。两类模型均通过统一接口向棒内气体组成、气腔压力及间隙传热模块传递裂变气体释放信息，从而实现裂变气体行为与燃料棒热--力响应的耦合。

\xsubsection{Forsberg--Massih裂变气体模型}{Forsberg--Massih Fission Gas Model}

裂变气体从燃料芯块释放至芯块--包壳间隙的主要机制包括反冲、击出与扩散。其中，反冲过程是指裂变发生位置距燃料表面小于裂变碎片在 $\mathrm{UO_2}$ 中平均自由程（约 $7\sim9~\upmu\mathrm{m}$）时，裂变产物直接逸出燃料表面的过程；击出过程则是裂变碎片穿过燃料表面时，将局部材料击出并伴随气体释放的现象；扩散过程则是裂变气体在浓度梯度驱动下，于晶粒内部迁移并向晶界输运的过程。

在不同温度条件下，上述机制的主导作用存在差异：在低温条件下，反冲与击出过程占主导；而在高温条件下，扩散过程成为裂变气体迁移与释放的主要机制。因此，在反应堆正常运行工况下，通常采用扩散模型对裂变气体行为进行描述。

为定量描述裂变气体在晶粒内部扩散及其在晶界处积累与释放过程，Forsberg 与 Massih\cite{forsbergDiffusionTheoryFission1985}在传统 Booth 模型基础上提出了一种改进模型。该模型将燃料晶粒理想化为半径为 $a$ 的等效球体，并假设裂变气体原子在晶粒内部均匀产生。在此假设下，晶粒内气体原子的扩散过程可由球坐标系下的扩散方程描述为：
\begin{equation}
	\frac{\partial C(r,t)}{\partial t}
	= D(t)\left(\frac{1}{r^2}\frac{\partial}{\partial r}\left(r^2 \frac{\partial C(r,t)}{\partial r}\right)\right)
	+ P(t)
\end{equation}
\noindent 式中：$C(r,t)$ 为半径位置 $r$ 处的裂变气体原子浓度（$\mathrm{m^{-3}}$），$D(t)$ 为温度相关的扩散系数（$\mathrm{m^2\,s^{-1}}$），$P(t)$ 为单位体积内气体的生成速率（$\mathrm{m^{-3}\,s^{-1}}$）。
与传统 Booth 模型中晶界被视为“完美吸收边界”的假设不同，Forsberg--Massih 模型考虑了晶界气体的有限吸收能力及再溶解效应。基于 Turnbull 提出的理论，其边界条件可表示为：
\begin{equation}
	C(a,t) = \frac{b(t)\,\delta\,N(t)}{2D(t)}
\end{equation}
\noindent 式中：$N(t)$ 表示单位晶界面积上的气体原子数密度（$\mathrm{m^{-2}}$），$b(t)$ 为再溶解速率常数（$\mathrm{s^{-1}}$），$\delta$ 为再溶解作用的特征深度（$\mathrm{m}$）。

为了耦合裂变气体模型和燃料性能分析程序，通常不是采用直接求解裂变气体扩散方程，而是采用解析解或者谱方法等数学变换，获得正交节点处的裂变气体行为。因此，接下来对裂变气体扩散方程进行基于核函数展开的降阶和卷积方程的状态变量化。

考虑扩散系数随温度变化的影响，引入时间变换变量：
\begin{equation}
	\tau = \int_0^t D(t') \, \mathrm{d}t'
\end{equation}

通过该变换，可将原扩散问题转化为等效常系数形式，从而获得解析表达。进一步推导可得晶界气体积累满足如下卷积形式：
\begin{equation}
	G_B(\tau) = \int_0^\tau K_2(\tau - \tau')\, q(\tau')\, \mathrm{d}\tau'
\end{equation}
\noindent 式中：$G_B$ 为晶界气体浓度，$K_2$ 为扩散核函数，$q(\tau)$ 为等效源项。

为提高数值计算效率，核函数采用指数形式近似为：
\begin{equation}
	1 - K_2(\tau) = \sum_{n=1}^{3} A_n \exp(-B_n \tau)
\end{equation}

在数值实现中，引入时间增量 $\Delta\tau = \tau^{n} - \tau^{n-1}$，并定义：
\begin{equation}
	f_n = \exp\left(-B_n \frac{\Delta\tau}{a^2}\right) - 1
\end{equation}

由此可建立晶粒内气体与晶界气体的递推更新关系，实现裂变气体扩散与积累过程的耦合求解。

在上述离散形式基础上，可进一步得到晶界气体累积过程的代数表达。根据 Forsberg--Massih 模型的解析解形式，对卷积积分进行离散化处理，可以得到解析解，其表达式为：
\begin{equation}
	a^2q =
	\frac{
		\frac{1}{2}\left(\beta(\tau_1)+\beta(\tau_2)\right)\Delta t
		- t_3
		+ h_4 \frac{\beta(\tau_2)}{D(\tau_2)} t_4
	}{
		- t_1 + \left(1 + h_4 \frac{\beta(\tau_2)}{D(\tau_2)}\right) t_2
	}
\end{equation}
\noindent 式中：各中间变量定义为：
\begin{align}
	t_1 & = \sum_{n=1}^{3} \frac{A_n}{B_n} f_n                                                            \\
	t_2 & = \frac{4}{\sqrt{\pi}}\left(\frac{\Delta\tau}{a^2}\right)^{3/2}
	- \frac{3}{2}\left(\frac{\Delta\tau}{a^2}\right)^2                                                    \\
	t_3 & = h_4 G_B(\tau_1)\left(\frac{\beta(\tau_2)}{D(\tau_2)} - \frac{\beta(\tau_1)}{D(\tau_1)}\right) \\
	t_4 & = \sum_{n=1}^{3} f_n G_n(\tau_1)
\end{align}

由上述离散表达式求得变量 $q$ 后，可进一步得到晶粒内部气体浓度$G_0$和晶界上的气体浓度$G_B$。更新形式为：
\begin{equation}
	G_B = \sum G_n \, G_n(\tau_2) = G_n(\tau_1) + f_n G_n(\tau_1) - a^2 q \frac{A_n}{B_n} f_n
\end{equation}

当晶界气体浓度达到饱和值时，将发生裂变气体释放，其饱和浓度可表示为：
\begin{equation}
	G_s = \frac{3}{2a} N_s
\end{equation}

当 $G_B > G_s$ 时，多余气体按一定比例释放，其释放量表示为：
\begin{equation}
	F_R = f \left(G_B - G_s\right)
\end{equation}

扩散系数 $D$ 是影响裂变气体迁移行为的关键参数，其大小主要受温度及燃耗等因素影响。在原始的Forsberg--Massih模型中，给出了基于实验数据拟合的经验模型。如下分段形式：
\begin{equation}
	D =
	\begin{cases}
		1.51\times10^{-17}\exp\left(-\dfrac{9508}{T}\right)                      & T < 1381~\mathrm{K}                     \\
		2.14\times10^{-13}\exp\left(-\dfrac{22884}{T}\right)\cdot f(\mathrm{Bu}) & 1381~\mathrm{K} \le T < 1650~\mathrm{K} \\
		1.09\times10^{-17}\exp\left(-\dfrac{6614}{T}\right)                      & T \ge 1650~\mathrm{K}
	\end{cases}
\end{equation}
\noindent 式中：$D$为裂变气体扩散系数（$\mathrm{m^2\,s^{-1}}$）；$T$为温度（$\mathrm{K}$）；$f(\mathrm{Bu})$为无量纲燃耗修正项，用于描述高燃耗条件下扩散行为的增强效应。

此外，在Trunbull等人\cite{denisSimulationIsothermalFission1996}的报告提出了经典的三项模型，包括高温热激活扩散、辐照增强扩散、裂变碎片直接导致的扩散。形式如下：
\begin{align}
	D_{\mathrm{Turnbull}} & = D_1 + D_2 + D_3\footnotemark                                                    \\
	D_1                   & = 7.6 \times 10^{-10} \exp\left(-\frac{35250}{T}\right)                           \\
	D_2                   & = 1.41 \times 10^{-25} \sqrt{\dot{F}} \exp\left(-\frac{13800}{T}\right) \cdot 4.0 \\
	D_3                   & = 2.0 \times 10^{-40} \dot{F}
\end{align}
\noindent 式中：$D_{\mathrm{Turnbull}}$、$D_1$、$D_2$和$D_3$均为扩散系数（$\mathrm{m^2\,s^{-1}}$）；$T$为温度，单位为$\mathrm{K}$；$\dot{F}$为体积裂变率（$\mathrm{m^{-3}\,s^{-1}}$）。
\footnotetext{在Bison报告中是直接线性相加，即1：1：1。但在原始的文献中采用的是1：4：4，即$D_{\mathrm{Turnbull}} = D_1 + 4D_2 + 4D_3$}

两个模型的扩散系数随温度变化的曲线如图\ref{fig:gas_diffusion}所示，两者都表现出扩散系数随温度上升而增加，但是Turnbull的模型还包含裂变碎片导致的扩散系数的增加。

为了实现上述的计算过程，本文基于时间步进策略构建如下的计算流程，通过MOOSE框架中的Material对象创建一个ForMas的材料类型，计算当前时刻晶界的裂变气体密度和释放量，各个状态变量采用时间步递推方式进行更新。

\begin{algorithm}[!htbp]
	\small
	\SetAlgoLined
	
	\caption{Forsberg--Massih裂变气体模型数值实现流程}
	\label{alg:forsberg_massih}
	
	\KwData{$\Delta t$, $T^n$, $P^n$,
		$G_i^{n-1}$, $G_B^{n-1}$, $\tau^{n-1}$}
	\KwResult{$G_i^n$, $G_B^n$, $\tau^n$, $F_R$}
	
	初始化 $G_i$、$G_B$ 和 $\tau$\;
	
	\For{每个时间步 $n$}{
		
		根据当前温度计算扩散系数 $D^n$\;
		
		根据当前裂变状态计算气体生成速率 $P^n$\;
		
		更新扩散时间变量
		$\tau^n=\tau^{n-1}+D^n\Delta t$\;
		
		计算扩散时间增量
		$\Delta\tau=\tau^n-\tau^{n-1}$\;
		
		\For{$i=1,2,3$}{
			计算核函数增量 $f_n^{(i)}$\;
		}
		
		计算核函数项 $t_1$ 和 $t_2$\;
		
		根据上一时间步晶界气体状态计算历史项
		$t_3$ 和 $t_4$\;
		
		求解辅助变量 $a_{2q}$\;
		
		\For{$i=1,2,3$}{
			更新晶内历史变量 $G_i^n$\;
		}
		
		更新晶界气体存量 $G_B^n$\;
		
		\eIf{$G_B^n>G_s$}{
			计算释放量 $F_R$\;
			
			更新释放后的晶界气体存量 $G_B^n$\;
		}{
			$F_R=0$\;
		}
	}
	
\end{algorithm}

\begin{figure}[!htb]
	\centering
	\begin{subfigure}[t]{0.45\textwidth}
		\centering
		\includegraphics[width=\linewidth,height=5.00cm,keepaspectratio]{gas_diffusion_forsberg.png}
		\caption{ForMas模型提供的扩散系数随温度变化曲线}
		\label{fig:gas_diffusion_forsberg}
	\end{subfigure}
	\hfill
	\begin{subfigure}[t]{0.45\textwidth}
		\centering
		\includegraphics[width=\linewidth,height=5.00cm,keepaspectratio]{gas_diffusion_turnbull.png}
		\caption{Turnbull修正之后的扩散系数曲线}
		\label{fig:gas_diffusion_turnbull}
	\end{subfigure}
	\caption{ForMas原始模型与Turnbull修正模型的裂变气体扩散系数对比}
	\label{fig:gas_diffusion}
\end{figure}

\xsubsection{细观机理裂变气体模型 SCIANTIX 及其接口实现}{Mechanistic Fission Gas Model SCIANTIX and Its Interface Implementation}
第二章已从燃料棒多物理场总体框架的角度，将SCIANTIX作为裂变气体行为的细观闭合模型，并给出了其与棒内气体状态及间隙传热之间的主要变量传递关系。本节在此基础上重点介绍SCIANTIX与Miracle程序的接口集成及状态变量映射方法，其内部裂变气体行为模型及理论基础可参见相关文献\cite{pizzocriSCIANTIXNewOpen2020}。

SCIANTIX采用C++开发，并具备作为独立程序运行或以模块形式与燃料性能程序耦合的能力。本文将SCIANTIX编译为静态库，并在Miracle的Make构建系统中配置相应的头文件搜索路径与库链接关系，从而实现SCIANTIX计算模块与Miracle主程序的静态链接。考虑到MOOSE框架中已定义核心基类\texttt{Material}，为避免与SCIANTIX源码中的同名类型发生命名冲突，本文将SCIANTIX中对应类重命名为\texttt{SciantixMaterial}，并对相关引用同步调整。

完成静态库集成后，本文基于MOOSE的\texttt{Material}基类开发封装类\texttt{SciantixWrapper}，作为Miracle与SCIANTIX之间的数据交换接口。该封装类在材料积分点尺度上读取温度、裂变率、应力等局部场变量，并将其转换为SCIANTIX所需的历史变量和模型输入；完成当前时间步的SCIANTIX计算后，再将裂变气体释放率、气体肿胀及其他内部状态量映射为MOOSE材料属性，实现裂变气体细观模型与燃料棒宏观多物理场计算之间的双向数据传递。

该封装类的核心功能如图\ref{fig:coupled_sciantix}所示，主要包括：首先，从 MOOSE 框架中读取温度、裂变率及应力等场变量，并完成单位转换；其次，将这些驱动变量按照SCIANTIX接口要求填充至历史变量数组中，同时设置模型选项及缩放参数；随后，在每个时间步内调用SCIANTIX求解器完成裂变气体行为的演化计算；最后，将计算得到的状态变量重新映射为 MOOSE 中的材料属性，如裂变气体释放率和气体肿胀等。详细的接口函数、状态变量及输入参数见~\ref{app:sciantix_interface}。

本文定义总的肿胀包含以下四部分，分别是晶内气泡肿胀、晶界气泡肿胀、晶内气体固溶肿胀和高燃耗肿胀。
\begin{equation}
	S_{\mathrm{gas}}
	=
	S_{\mathrm{intra}}
	+
	S_{\mathrm{inter}}
	+
	S_{\mathrm{sol}}
	+
	f_{\mathrm{HBS}} \, P_{\mathrm{HBS}}
\end{equation}
\begin{figure}[!htb]
	\centering
	\includegraphics[width=15.5cm,height=11.00cm,keepaspectratio]{coupled_sciantix.pdf}
	\caption{Miracle与SCIANTIX的耦合接口及数据传递关系}
	\label{fig:coupled_sciantix}
\end{figure}

通过上述方式，实现了SCIANTIX与MOOSE 框架之间的数据双向传递与状态变量的持续演化，从而保证了裂变气体行为模型在多物理场耦合计算中的一致性与稳定性。

\xsubsection{间隙传热模型实现}{Gap Heat Transfer Model Implementation}

第二章已给出燃料--包壳间隙传热的总体关系及其与燃料棒热--力--气体行为之间的耦合框架，因此本节不再重复相关基本关系，而重点介绍间隙传热模型的具体实现及其与裂变气体释放行为的数据耦合方式。间隙传热状态直接影响燃料与包壳之间的热量传递，并通过温度变化进一步影响热膨胀、间隙宽度及接触状态；与此同时，燃料中释放的Xe和Kr进入棒内自由体积后，会改变间隙气体组成、气腔压力及混合气体热物性，从而反馈影响气体传导传热。因此，在程序实现中需要统一处理几何间隙与接触状态、混合气体传导、固--固接触传导、热辐射以及裂变气体释放引起的气体状态变化。

基于MOOSE框架的模块化特性，本文采用Action-Object分离的设计思想，如图\ref{fig:thermal_contact}所示，实现了间隙传热模型的自动构建与物理过程解耦。其中，整体实现由三个核心组件构成：\texttt{LWRThermalContactAction}、\texttt{GasGapConductance}以及\texttt{GasMixture}。

首先，\texttt{LWRThermalContactAction}作为入口模块，负责统一创建与间隙传热相关的所有对象。该Action在不同的执行阶段（如\texttt{add\_aux\_variable}、\texttt{add\_aux\_kernel}、\texttt{add\_material}和\texttt{add\_bc}）动态构建计算所需的变量与算子。具体而言，其核心功能包括：

（1）构建间隙几何相关变量。通过\texttt{PenetrationAux}与\texttt{GapValueAux}分别计算接触面之间的法向距离以及跨界面的温度映射，从而实现主从表面之间的信息传递。

（2）实现固体导热率的跨界面映射。首先利用\texttt{MaterialRealAux}将块内定义的导热率提取为全局辅助变量，然后通过\texttt{GapValueAux}将主侧导热率映射至从侧，实现双侧材料属性在间隙模型中的统一调用。

（3）自动生成边界条件。根据用户指定的导热模型类型（如默认\texttt{GapPerfectConductance}或者其他已知的气体导热模型），在接触边界上构建相应的热通量边界条件，实现温度场的耦合求解。

在此基础上，\texttt{GasGapConductance}作为Material对象，负责在各积分点计算总间隙传热系数。该模型通过\texttt{computeQpConductance()}将间隙传热分解为气体传导、热辐射和固--固接触传导三部分：
\begin{equation}
	h_{\text{gap}}
	=
	h_{\text{cond}}
	+
	h_{\text{rad}}
	+
	h_{\text{contact}}
\end{equation}
\noindent 式中：$h_{\text{cond}}$为气体传导传热系数（$\mathrm{W\,m^{-2}\,K^{-1}}$），主要受混合气体热导率、间隙尺寸、气体压力及温度跳跃距离影响；$h_{\text{rad}}$为辐射传热系数（$\mathrm{W\,m^{-2}\,K^{-1}}$）；$h_{\text{contact}}$为固--固接触传导传热系数（$\mathrm{W\,m^{-2}\,K^{-1}}$），与接触压力、表面粗糙度及两侧材料热物性有关。上述物理关系及参数模型已在第2章给出总体形式，其中等效间隙长度与几何修正见\ref{app:gap_geometry_models}，Lanning--Hann和Toptan温度跳跃模型见\ref{app:gap_temperature_jump_models}，混合气体热导率模型见\ref{app:gap_default_gas_conductivity}，固--固接触传导补充关系见\ref{app:gap_contact_supplementary_models}。

在气体传导建模中，程序提供多种温度跳跃模型，用于考虑稀薄气体条件下分子自由程对界面传热的影响。对于高温或较宽压力范围内的气体热导率计算，还可采用\ref{app:gap_advanced_gas_conductivity}所述温度--压力相关模型。此外，该Material通过耦合变量\texttt{neighbor\_conductivity}与\texttt{secondary\_conductivity}获取间隙两侧固体材料的热导率，从而将燃料与包壳热物性共同引入间隙传热计算。

为了描述气体成分随时间的演化过程，引入了\texttt{GasMixture}向量后处理器。该模块以后处理器形式运行，不参与残差计算，而是用于提供时间依赖的气体组成信息。其主要功能包括：

（1）初始化气体组成，根据初始气体类型及其摩尔分数计算各组分初始摩尔数；

（2）处理裂变气体释放，通过外部postprocessor输入释放量，并按设定比例分配至不同气体组分；

（3）模拟燃料棒再制造过程，包括气体重置或持续释放两种模式；

（4）输出各气体组分随时间变化的摩尔数，并通过\texttt{VectorPostprocessor}接口提供给\texttt{GasGapConductance}调用。

在\texttt{GasGapConductance}中，通过读取\texttt{GasMixture}提供的气体摩尔数向量，计算各组分摩尔分数，并进一步用于混合气体热导率、热适应系数以及跳跃距离的计算。这种“后处理器 + 材料模型”的耦合方式避免了在主方程中显式引入额外自由度，从而提高了计算效率。

整体而言，该实现通过Action自动构建计算结构，Material负责物理模型计算，VectorPostprocessor提供间隙内气体成分信息，实现了间隙传热、气体行为与接触几何之间的紧密耦合。

\begin{figure}[!htb]
	\centering
	\includegraphics[width=15.5cm,height=11.00cm,keepaspectratio]{thermal_contact.pdf}
	\caption{LWRThermalContact间隙传热模型的计算流程}
	\label{fig:thermal_contact}
\end{figure}

\xsubsection{裂变气体释放与间隙传热的耦合关系}{Coupling between Fission Gas Release and Gap Heat Transfer}
如前所述，本文间隙传热模型显式考虑不同气体组分对混合气体热物性的影响，其中包括初始填充气体以及由燃料释放的Xe和Kr等裂变气体。随着燃耗和燃料温度的演化，裂变气体持续由燃料内部释放至棒内自由体积，使气体组分随辐照历史动态变化。由于不同气体组分具有不同的热导率，Xe和Kr含量的变化将进一步改变混合气体的等效热导率，并由此影响燃料--包壳间隙的气体传导能力。因此，在整个燃料棒服役过程中，间隙气体热物性并非固定参数，而是与裂变气体释放行为和棒内气体状态动态耦合。

将前述裂变气体模型与间隙传热模型进行耦合后的整体流程如图\ref{fig:coupled_sciantix_with_gap}所示。在当前时间步内，首先由裂变气体释放模型（ForMas或SCIANTIX）计算时刻 $t^{n}$ 的裂变气体释放量及其组分；随后通过 VectorPostprocessor 类 \texttt{GasMixture} 更新间隙内各气体组分的物质的量与摩尔分数；在此基础上，\texttt{GasGapConductance} 根据用户选定的导热模型分别计算跳跃距离、各组分热导率以及混合气体热导率；之后，GapHeatTransfer 利用更新后的间隙传热系数求解下一时刻 $t^{n+1}$ 的温度分布；最终，新的温度场再反馈至裂变气体模块，用于计算时刻 $t^{n+1}$ 的裂变气体演化行为。通过上述迭代过程，可实现裂变气体释放与间隙传热之间的双向耦合。

\begin{figure}[!htb]
	\centering
	\includegraphics[width=13.5cm,height=9.00cm,keepaspectratio]{coupled_sciantix_with_gap.pdf}
	\caption{裂变气体释放、气体组分与间隙传热系数的耦合关系}
	\label{fig:coupled_sciantix_with_gap}
\end{figure}

\xsection{程序验证与确认}{Verification and Validation of the Implementation}
为验证本章所实现的非线性力学本构积分、裂变气体行为模型以及间隙传热耦合模块的正确性，本文按照“局部本构验证—子模型验证—耦合接口验证”的思路设计了一系列算例。对于力学非线性问题，首先采用单元级算例考察径向返回算法在蠕变和硬化模型中的应力更新能力，并将 UO$_2$、MOX 及典型包壳材料的数值结果与理论参考解或相应基准关系进行对比。对于裂变气体行为，分别对简化机理模型 ForMas 和 SCIANTIX 耦合接口进行验证：前者通过晶界气体累积与释放算例检验模型对释放阈值、扩散系数和释放比例的响应，后者采用 SCIANTIX 回归测试中的代表性算例，比较 Miracle 耦合计算与基准结果在晶内气泡、晶界释放、高燃耗结构和化学计量偏差等变量上的一致性。对于间隙传热模型，则构建二维接触传热算例，比较不同气体热导率模型和压力反馈条件下间隙换热系数的变化规律，以验证气体温度、压力和间隙宽度等变量的耦合传递是否合理。

\xsubsection{径向返回验证}{Verification of the Radial Return Algorithm}
\subsubsection{蠕变}
根据总蠕变$\epsilon_{\mathrm{creep}}$由各种蠕变组成：
\begin{equation}
	\epsilon_{\mathrm{creep}} = \epsilon_{thermal\, creep} + \epsilon_{irradiation\, creep}
\end{equation}

依据第4.2.2节介绍的径向返回算法，本文构建了不同的蠕变模型，包括热蠕变和辐照蠕变。具体的模型见表\ref{tab:creep_classes}，不同材料采用独立的蠕变应力更新类进行描述，各类均继承自 MOOSE 的应力更新框架，并在径向返回算法中实现蠕变应变的演化计算。相关关系如图\ref{fig:creep_flow_chat}所示。
\begin{table}[H]
	\centering
	\caption{与蠕变行为相关的主要类及其功能}
	\label{tab:creep_classes}
	\begin{tabularx}{\textwidth}{@{}
		>{\centering\arraybackslash}m{2.2cm}
		>{\centering\arraybackslash}m{3.4cm}
		>{\centering\arraybackslash}X
		@{}}
		\toprule
		\multicolumn{1}{c}{材料类型} & \multicolumn{1}{c}{类名}                                     & {\centering 功能描述\par} \\
		\midrule

		D9
		                         & D9CreepStressUpdate
		                         & 实现 D9 合金的蠕变应力更新算法，用于描述高温条件下材料的蠕变变形行为                                           \\

		HT9
		                         & HT9\allowbreak Creep\allowbreak Stress\allowbreak Update
		                         & 实现 HT9 钢的蠕变本构模型，考虑温度和应力对蠕变速率的影响                                                    \\

		SS316
		                         & SS316\allowbreak Creep\allowbreak Stress\allowbreak Update
		                         & 实现 SS316 不锈钢的蠕变演化模型，用于模拟长期服役过程中的应变积累                                               \\

		MOX
		                         & MOXCreepStress
		                         & 实现MOX的蠕变演化模型，考虑温度、应力和裂变率对蠕变速率的影响                                                   \\

		$UO_2$
		                         & UO2CreepStress
		                         & 实现$UO_2$的蠕变演化模型，考虑温度、应力和裂变率对蠕变速率的影响                                                \\

		\bottomrule
	\end{tabularx}
\end{table}

\begin{figure}[!htb]
	\centering
	\includegraphics[width=13.5cm,height=9.00cm,keepaspectratio]{creep_flow_chat.pdf}
	\caption{基于径向返回算法的蠕变模型类继承关系图}
	\label{fig:creep_flow_chat}
\end{figure}

针对蠕变模型验证，采用图\ref{fig:case} 所示的$1\times1\times1$立方体单元，在指定加载面施加大小为$10\,\mathrm{MPa}$的恒定法向面力，并在三个坐标方向施加对称边界条件，以构建单轴应力状态。计算过程中温度保持恒定，为 $1000\,\mathrm{K}$。整个计算共划分为 10 个时间步，每个时间步长为 $1\times10^4\,\mathrm{s}$，总模拟时间为 $1\times10^5\,\mathrm{s}$。

\begin{figure}[!htb]
	\centering
	\includegraphics[width=9.00cm,height=6.75cm,keepaspectratio]{creep_case.pdf}
	\caption{蠕变模型使用的$1\times1\times1$ 立方体单元验证模型示意图}
	\label{fig:case}
\end{figure}

为构造独立于程序数值积分过程的理论参考解，本算例采用恒定温度和恒定单轴载荷条件。对于给定的材料蠕变本构关系：
\begin{equation}
	\dot{\bar{\varepsilon}}_{\mathrm{cr}}
	=
	f\left(
	T,\bar{\sigma},\varPhi,\dot{\varPhi},t
	\right)
\end{equation}

在$T=T_0$、$\bar{\sigma}=\bar{\sigma}_0$以及给定中子通量条件下，理论蠕变应变可由本构关系直接积分得到：
\begin{equation}
	\bar{\varepsilon}_{\mathrm{cr}}^{\mathrm{ref}}(t)
	=
	\int_0^t
	f\left(
	T_0,\bar{\sigma}_0,
	\varPhi(\tau),\dot{\varPhi},\tau
	\right)
	\,\mathrm{d}\tau
	\label{eq:creep_reference_solution}
\end{equation}

对于本构关系已经以应变--时间显式形式给出的模型，则直接代入相应温度、应力和时间即可获得理论参考值；对于以蠕变率形式给出的模型，则按照式~\eqref{eq:creep_reference_solution}进行时间积分。该理论参考解不经过有限元本构积分过程，因此可用于独立检验程序中径向返回算法及材料状态更新的实现正确性。

图~\ref{fig:ht9_creep_verification}--图~\ref{fig:d9_creep}给出了不同包壳材料蠕变模型的程序数值结果与由相应本构关系直接积分得到的理论参考解之间的对比。对于HT9材料，分别测试了是否考虑一阶蠕变、采用MFH默认二阶蠕变模型或RKY二阶蠕变模型、是否引入三阶蠕变项以及是否考虑辐照蠕变等不同组合条件下的模型响应。对于SS316和D9材料，则重点验证热蠕变与辐照蠕变模型在长期加载条件下的应变演化行为。

结果表明，各材料蠕变模型的程序数值结果与理论参考解保持良好一致，说明本文所实现的蠕变本构积分及状态更新过程能够正确再现相应本构关系所规定的时间相关变形行为。

\begin{figure}[!htb]
	\centering
	\begin{subfigure}[t]{0.45\textwidth}
		\centering
		\includegraphics[width=\linewidth,height=6.7500cm,keepaspectratio]{ht9_creep_MFH.png}
		\caption{MFH 蠕变模型验证}
		\label{fig:ht9_creep_MFH}
	\end{subfigure}
	\hfill
	\begin{subfigure}[t]{0.45\textwidth}
		\centering
		\includegraphics[width=\linewidth,height=5.00cm,keepaspectratio]{ht9_creep_RKY.png}
		\caption{RKY 蠕变模型验证}
		\label{fig:ht9_creep_RKY}
	\end{subfigure}
	\caption{HT9蠕变模型数值结果与理论参考解对比验证}
	\label{fig:ht9_creep_verification}
\end{figure}
\begin{figure}[!htb]
	\centering
	\begin{subfigure}[t]{0.45\textwidth}
		\centering
		\includegraphics[width=\linewidth,height=5.00cm,keepaspectratio]{ss316_creep_complex.png}
		\caption{Complex 辐照蠕变和热蠕变数据验证}
		\label{fig:ss316_creep_complex}
	\end{subfigure}
	\hfill
	\begin{subfigure}[t]{0.45\textwidth}
		\centering
		\includegraphics[width=\linewidth,height=5.00cm,keepaspectratio]{ss316_creep_bison.png}
		\caption{BISON辐照蠕变和热蠕变模型验证}
		\label{fig:ss316_creep_bison}
	\end{subfigure}
	\caption{SS316辐照蠕变模型数值结果与理论参考解对比验证}
	\label{fig:ss316_creep_verification}
\end{figure}
\begin{figure}[!htb]
	\centering
	\includegraphics[width=9.00cm,height=6.75cm,keepaspectratio]{d9_creep.png}
	\caption{D9辐照蠕变与热蠕变模型数值结果和理论参考解对比}
	\label{fig:d9_creep}
\end{figure}
在燃料蠕变模型验证中，本文分别选取UO$_2$和MOX燃料开展单元立方体算例。相关蠕变本构关系已在前文给出，其中MOX燃料的瞬态蠕变和稳态蠕变关系分别见式\eqref{eq:mox_transient_creep}和式\eqref{eq:mox_steady_creep}，UO$_2$燃料的MATPRO蠕变模型及MOX蠕变模型的补充参数见\ref{app:mox_mechanical_correlations}。因此，本节不再重复列出模型公式，而重点验证上述本构关系在材料积分及全局力学求解过程中的数值实现正确性。

验证算例采用$1\times1\times1$单元立方体模型，在$y$方向施加单轴载荷或位移边界条件，并给定相应的恒定温度、裂变率及时间步长。对于给定载荷历史，理论参考结果由相应蠕变本构关系直接积分获得；程序计算结果则由有限元求解得到的轴向位移或等效蠕变应变提取。通过二者对比，可检验蠕变应变更新、材料状态变量演化以及局部本构积分与全局力学求解之间的数据传递是否正确。

对于UO$_2$燃料，轴向总位移的程序计算结果为$8.9598\times10^{-4}\,\mathrm{m}$，理论参考结果为$9.07136\times10^{-4}\,\mathrm{m}$，相对差异约为$1.23\%$。对于MOX燃料，本文采用包含加载、应力松弛和二次加载阶段的单元算例，以同时检验瞬态蠕变和稳态蠕变项的数值实现。第一阶段理论参考蠕变应变为$9.343657\times10^{-3}$，程序计算结果为$9.4418\times10^{-3}$，相对差异约为$1.05\%$。上述结果表明，UO$_2$和MOX燃料蠕变模型在当前程序框架下均能够得到与相应理论参考结果一致的数值响应，说明相关蠕变本构积分及其与全局力学求解的耦合实现正确。

\subsubsection{硬化}
包壳材料在堆内服役过程中会受到温度、快中子辐照以及机械载荷的共同作用，其屈服强度和塑性变形能力均会随服役状态发生变化。对于燃料性能分析程序而言，硬化模型的作用在于描述塑性应变累积后屈服面的扩张规律，并为径向返回算法提供屈服应力及其对内变量的导数。因此，硬化模型不仅影响局部应力更新的准确性，也会影响全局非线性迭代的收敛稳定性。

不同包壳材料的极限拉伸强度$\sigma_{\mathrm{uts}}$和最小延伸率$\epsilon_t$关系及相关参数见~\ref{app:cladding_supplementary_models}。参考燃料性能分析程序中常用的工程化处理方式，本文首先采用各向同性线性硬化模型对包壳材料的塑性应变硬化行为进行近似描述。考虑到部分材料缺少完整的真应力--塑性应变曲线，本文采用基于屈服强度、抗拉强度和延伸率的工程近似关系估算线性硬化参数：
\begin{equation}
	H = \frac{\sigma_{\mathrm{uts}}-\sigma_y}{\epsilon_t}
	\label{eq:linear_hardening_modulus}
\end{equation}
\noindent 式中： $\sigma_y$ 为初始屈服强度（$\mathrm{Pa}$）。该模型采用线性关系近似描述屈服应力随累积等效塑性应变的变化。由于抗拉强度和延伸率并不直接构成材料的真实应力--塑性应变曲线，式中$H$属于工程等效硬化参数，而非由完整塑性曲线确定的严格切线硬化模量；其适用范围受材料状态及试验数据条件限制。

在程序实现方面，本文针对 D9、HT9 和 SS316 等包壳材料分别实现了各向同性硬化应力更新模型。各模型均嵌入 MOOSE 框架的径向返回积分流程中，并通过统一的塑性应力更新接口与弹塑性本构积分算法耦合。对于给定的总应变增量，程序首先进行弹性预测并判断试应力状态是否超过屈服面；若发生屈服，则根据硬化模型返回当前屈服应力、硬化变量及其导数，并通过局部牛顿迭代求解塑性乘子增量。由此，硬化模型可作为材料相关模块独立提供屈服面演化规律，而不改变径向返回算法的总体结构。

为了验证各硬化模型的实现正确性，本文采用如图 \ref{fig:hardening_case} 所示的 $1\times1\times1$ 立方体单元开展单轴拉伸数值算例。计算中，立方体底端施加固定约束，顶端施加恒定位移加载速率 $1\times10^{-8}(m/s)$，其余表面采用自由边界条件，模拟总时间为 $2\times10^6$ s。该算例能够在简单应力状态下产生单调增加的等效塑性应变，从而便于将程序输出的硬化后屈服应力与理论硬化关系进行逐点对比。

不同包壳材料对应的硬化模型类及其主要功能如表 \ref{tab:hardening_classes} 所示。上述模型均继承自统一的塑性应力更新框架，并通过重载屈服应力、硬化变量及其导数等关键函数，实现材料参数、硬化规律和温度相关性的差异化描述。图 \ref{fig:hardening_flow_chat} 给出了各向同性硬化模型在 MOOSE 框架中的类层次结构及继承关系。
\begin{figure}[!htb]
	\centering
	\includegraphics[width=9.00cm,height=6.75cm,keepaspectratio]{hardening_case.pdf}
	\caption{硬化模型验证采用的 $1\times1\times1$ 立方体单元示意图}
	\label{fig:hardening_case}
\end{figure}

\begin{table}[!htbp]
\centering
\caption{与各向同性硬化行为相关的主要类及其功能}
\label{tab:hardening_classes}
\begin{tabularx}{\textwidth}{@{}
  >{\centering\arraybackslash}m{2.0cm}
  >{\centering\arraybackslash}m{5.2cm}
  >{\centering\arraybackslash}X
  @{}}
\toprule
材料类型 & 类名 & 功能描述 \\
\midrule

	D9
	& Isotropic\allowbreak D9\allowbreak Hardening\allowbreak Stress\allowbreak Update
	& 实现D9合金的各向同性硬化模型，用于描述塑性变形过程中屈服面的均匀扩张行为，并考虑温度对硬化参数的影响 \\

	HT9
	& Isotropic\allowbreak HT9\allowbreak Hardening\allowbreak Stress\allowbreak Update
	& 实现HT9铁素体/马氏体钢的各向同性硬化本构模型，用于描述塑性应变累积导致的材料强化效应 \\

	SS316
	& Isotropic\allowbreak SS316\allowbreak Hardening\allowbreak Stress\allowbreak Update
	& 实现SS316不锈钢的各向同性硬化模型，支持Ludwik、Chaboche以及Armstrong等多种硬化规律，并可考虑温度、晶粒尺寸及厚度等因素对硬化行为的影响 \\

\bottomrule
\end{tabularx}
\end{table}

\begin{figure}[!htb]
	\centering
	\includegraphics[width=11.5cm,height=7.00cm,keepaspectratio]{hardening_flow_chat.pdf}
	\caption{基于径向返回算法的硬化模型类继承关系图}
	\label{fig:hardening_flow_chat}
\end{figure}

图 \ref{fig:hardening_verification} 给出了 D9、HT9 和 SS316 材料在单轴加载过程中硬化后屈服应力随等效塑性应变的变化关系，并将数值结果与相应理论关系进行对比。对于 D9 和 HT9，本文采用线性各向同性硬化模型，因此屈服应力随等效塑性应变近似呈线性增长。对于 SS316，程序中实现了包括线性硬化和幂律硬化在内的多种硬化关系，以覆盖不同材料关联式下的强化行为。可以看出，各材料模型的数值计算结果均能够与理论曲线保持良好一致，说明程序能够正确传递等效塑性应变、更新屈服应力，并在径向返回过程中保持硬化变量的一致性。

\begin{figure}[!htb]
	\centering
	\begin{subfigure}[b]{0.45\textwidth}
		\centering
		\includegraphics[width=\linewidth,height=5.00cm,keepaspectratio]{D9_hardening_verify.png}
		\caption{D9 硬化模型验证}
		\label{fig:D9_hardening_verify}
	\end{subfigure}
	\hfill
	\begin{subfigure}[b]{0.45\textwidth}
		\centering
		\includegraphics[width=\linewidth,height=5.00cm,keepaspectratio]{HT9_hardening_verify.png}
		\caption{HT9 硬化模型验证}
		\label{fig:HT9_hardening_verify}
	\end{subfigure}
	\vspace{0.4cm}
	\begin{subfigure}[b]{0.45\textwidth}
		\centering
		\includegraphics[width=\linewidth,height=5.00cm,keepaspectratio]{SS316_hardening_verify.png}
		\caption{SS316 硬化模型验证}
		\label{fig:SS316_hardening_verify}
	\end{subfigure}

	\caption{不同包壳材料硬化本构模型的验证结果}
	\label{fig:hardening_verification}
\end{figure}

\xsubsection{裂变气体模型的数值测试与基准对比}{Numerical Testing and Benchmark Comparison of Fission Gas Models}
裂变气体行为会通过燃料棒内压、燃料肿胀以及燃料--包壳间隙传热条件影响燃料元件的整体性能。为验证本文实现的裂变气体模型在气体产生、晶界积累和释放判据等关键环节上的正确性，本节不以实验数据对模型物理预测能力进行确认，而是针对ForMas模型的数值实现以及SCIANTIX外部程序耦合接口开展基准测试与一致性检查。ForMas部分主要考察气体产生、晶界积累、饱和释放及关键参数响应是否符合模型理论预期；SCIANTIX部分则通过其回归测试算例检查输入输出映射、状态变量更新、时间推进及结果回写过程是否与原程序保持一致。裂变气体模型面向辐照实验数据的定量评价将在第五章棒级验证中进一步开展。

\subsubsection{ForMas模型基准测试}
裂变气体简化机理模型 ForMas 的验证采用文献\cite{forsbergDiffusionTheoryFission1985}给出的基准算例。该算例通过给定温度历程和裂变气体产生率，考察晶内扩散、晶界气体累积以及饱和释放过程的时序演化。计算中所采用的温度历程如图\ref{fig:ForMas_temperature_history}所示，裂变气体产生率、线功率与温度之间的对应关系见表\ref{tab:ForMas_condition}。

\begin{figure}[!htb]
	\centering
	\includegraphics[width=9.00cm,height=6.75cm,keepaspectratio]{ForMas_temperature_history.pdf}
	\caption{ForMas基准算例的温度历程}
	\label{fig:ForMas_temperature_history}
\end{figure}

\begin{table}[htbp]
	\centering
	\caption{温度、线功率和裂变气体产生率的对应关系（数据来源：文献\cite{forsbergDiffusionTheoryFission1985}）}
	\label{tab:ForMas_condition}

	\begin{tabular*}{\textwidth}{@{\extracolsep{\fill}}cccc@{}}
		\toprule

		温度/${}^\circ\mathrm{C}$ &
		线功率/$\mathrm{kW\,m^{-1}}$ &
		比功率/$\mathrm{kW\,kg^{-1}}$ &
		气体生成率/$\mathrm{mol\,m^{-3}\,s^{-1}}$ \\

		\midrule

		1000 & 26.8 & 29.8 & \(4.4 \times 10^{-6}\) \\
		1100 & 30.6 & 34.1 & \(5.1 \times 10^{-6}\) \\
		1200 & 34.4 & 38.3 & \(5.7 \times 10^{-6}\) \\
		1300 & 38.2 & 42.4 & \(6.4 \times 10^{-6}\) \\
		1400 & 41.9 & 46.6 & \(7.2 \times 10^{-6}\) \\
		1500 & 45.7 & 50.8 & \(7.7 \times 10^{-6}\) \\
		1600 & 49.4 & 55.0 & \(8.3 \times 10^{-6}\) \\
		1700 & 53.2 & 59.2 & \(8.9 \times 10^{-6}\) \\
		1800 & 57.0 & 63.4 & \(9.6 \times 10^{-6}\) \\
		1900 & 61.0 & 67.6 & \(1.0 \times 10^{-5}\) \\
		2000 & 64.6 & 71.8 & \(1.1 \times 10^{-5}\) \\

		\bottomrule
	\end{tabular*}
\end{table}

图\ref{fig:Gb_and_Gr}给出了基准工况下晶界气体覆盖度和累计释放密度的演化结果。可以看出，晶界气体密度与饱和气体密度之比随时间逐渐升高；当其接近饱和阈值时，模型触发晶界气体释放，随后晶界气体重新进入积累阶段，并在再次达到饱和条件后发生下一次释放。因此，覆盖度曲线呈现周期性累积--释放特征。与此对应，累计释放密度曲线表现为阶跃式增长，其中平台段代表两次释放事件之间的气体积累过程。该结果表明，本文实现的 ForMas 模型能够正确描述晶界饱和判据、瞬时释放以及释放后重新积累的基本物理过程，其演化规律与已有研究结果一致\cite{ZhangXiaoDiHeRanLiaoXinKuaiLieBianQiTiShiFangXingWeiMoNiYanJiu2016}。

为进一步检验模型对关键参数的响应，本文比较了不同扩散系数模型和单次释放比例 $f$ 条件下的计算结果，如图\ref{fig:Gb_and_Gr_compare}所示。在相同热工况下，Turnbull 扩散模型对应的释放次数明显多于原始 ForMas 扩散模型。结合扩散系数随温度变化的关系可知，Turnbull 模型在相同温度下给出的扩散系数较大，使裂变气体更快向晶界迁移并达到饱和条件，因此更频繁地触发释放事件，累计释放量也相对较高。相比之下，单次释放比例 $f$ 对最终累计释放量的影响较弱，主要改变释放事件的频率和单次释放幅度。这是由于给定时间区间内裂变气体总生成量和晶界饱和阈值基本保持一致；当 $f$ 较小时，单次释放量减少，但释放后晶界气体更容易在后续阶段重新累积至饱和状态。上述参数响应符合 ForMas 模型的物理预期，说明本文实现能够合理反映扩散系数和释放比例对裂变气体释放行为的影响。
\begin{figure}[!htb]
	\centering
	\begin{subfigure}[t]{0.45\textwidth}
		\centering
		\includegraphics[width=\linewidth,height=5.00cm,keepaspectratio]{Gb_Gs.pdf}
		\caption{晶界气体覆盖度随时间变化}
		\label{fig:Gb_Gs}
	\end{subfigure}
	\hfill
	\begin{subfigure}[t]{0.45\textwidth}
		\centering
		\includegraphics[width=\linewidth,height=5.00cm,keepaspectratio]{gas_released_density.pdf}
		\caption{裂变气体释放密度随时间变化}
		\label{fig:gas_released_density}
	\end{subfigure}
	\caption{晶界气体积累与裂变气体释放行为随时间的演化}
	\label{fig:Gb_and_Gr}
\end{figure}

\begin{figure}[!htb]
	\centering
	\begin{subfigure}[t]{0.45\textwidth}
		\centering
		\includegraphics[width=\linewidth,height=5.00cm,keepaspectratio]{Gb_Gs_compare.pdf}
		\caption{不同扩散模型和释放率下的晶界气体覆盖度随时间变化}
		\label{fig:Gb_Gs_compare}
	\end{subfigure}
	\hfill
	\begin{subfigure}[t]{0.45\textwidth}
		\centering
		\includegraphics[width=\linewidth,height=5.00cm,keepaspectratio]{gas_released_density_compare.pdf}
		\caption{不同扩散模型和释放率下的裂变气体释放密度随时间变化}
		\label{fig:gas_released_density_compare}
	\end{subfigure}
	\caption{不同扩散模型和释放率条件下裂变气体行为演化的比较}
	\label{fig:Gb_and_Gr_compare}
\end{figure}

\subsubsection{SCIANTIX耦合接口一致性测试}
SCIANTIX是面向核燃料微观行为分析的机理程序，已针对裂变气体行为、气泡演化、高燃耗结构以及燃料化学状态等问题开展了系统验证\cite{giozuSciantixSciantixofficialSCIANTIX2025,zulloSCIANTIXCodeFission2023,pizzocriSCIANTIXNewOpen2020}。本文并不重新验证 SCIANTIX 原程序本身，而是重点验证其与本文燃料性能分析程序耦合后的调用流程、输入输出映射和状态变量更新是否正确。为此，本文从 SCIANTIX 回归测试体系中选取四类具有代表性的算例，分别覆盖晶内气泡演化、晶界气体释放、高燃耗结构演化以及燃料氧化状态变化等过程。各算例所对应的主要物理领域、验证变量和验证目的见表\ref{tab:verification}。

\begin{table}[htbp]
	\centering
	\caption{SCIANTIX耦合验证算例及核心验证变量}
	\label{tab:verification}

	\begin{tabularx}{\textwidth}{@{}
		>{\centering\arraybackslash}m{2.2cm}
		>{\centering\arraybackslash}m{1.6cm}
		>{\centering\arraybackslash}m{3.3cm}
		>{\centering\arraybackslash}X
		@{}}
		\toprule
		物理领域 & 测试模块                    & 核心验证变量 & 验证目的 \\
		\midrule

		晶内物理
		     & Baker
		     & 晶内气泡肿胀
		     & 验证气泡成核、生长及晶内气体肿胀行为                      \\

		晶界与释放
		     & White
		     & 晶界气体肿胀、释放分数
		     & 验证晶界气泡演化、饱和释放及burst释放响应                 \\

		组织演化
		     & HBS
		     & 高燃耗结构孔隙率
		     & 验证高燃耗条件下组织重构及孔隙形成过程                     \\

		热化学
		     & Oxidation
		     & 化学计量偏差
		     & 验证燃料氧化状态及其对相关状态变量的影响                    \\

		\bottomrule
	\end{tabularx}
\end{table}

验证结果如图\ref{fig:bark}--图\ref{fig:academic_deviation}所示。Baker 算例用于检验晶内气泡演化过程，图\ref{fig:bark}表明，耦合程序得到的裂变气体释放分数和晶内气泡肿胀随时间变化趋势与 SCIANTIX 基准结果保持一致，说明晶内气体生成、气泡生长及相关状态变量反馈过程能够被正确传递。White 算例进一步考察晶界气泡演化和释放行为，图\ref{fig:white}显示，在包含 burst 释放特征的工况下，耦合计算仍能捕捉裂变气体释放分数和晶界气体肿胀的主要变化特征，表明晶界饱和、释放触发及释放后状态更新过程实现正确。

对于高燃耗结构算例，图\ref{fig:hbs}给出了裂变气体释放分数和气体肿胀的对比结果。计算结果能够反映高燃耗阶段组织重构后孔隙形成与气体行为之间的耦合关系，说明 SCIANTIX 中与 HBS 相关的状态变量能够通过接口在本文程序中稳定更新。图\ref{fig:academic_deviation}给出了燃料化学计量偏差的演化结果。可以看出，化学计量偏差对时间步长较为敏感；当时间步长过大时，快速变化阶段的响应会出现一定偏差，而减小时间步长后数值结果能够明显改善。这表明该过程在耦合计算中需要采用足够细的时间离散，以保证热化学状态变量更新的精度。

综上，四类代表性算例从晶内、晶界、组织演化和热化学行为等不同层面对 SCIANTIX 耦合接口进行了验证。结果表明，本文建立的耦合程序能够正确完成工况输入、状态变量传递、外部程序调用和结果回写，并能够保持与 SCIANTIX 基准结果一致的主要物理演化趋势。
\begin{figure}[!htb]
	\centering
	\begin{subfigure}[t]{0.45\textwidth}
		\centering
		\includegraphics[width=\linewidth,height=5.00cm,keepaspectratio]{academic_fgr_bark.png}
		\caption{SCIANTIX 与 Miracle 裂变气体释放分数对比}
		\label{fig:academic_fgr_bark}
	\end{subfigure}
	\hfill
	\begin{subfigure}[t]{0.45\textwidth}
		\centering
		\includegraphics[width=\linewidth,height=5.00cm,keepaspectratio]{academic_swelling_bark.png}
		\caption{SCIANTIX 与 Miracle 晶内气泡肿胀对比}
		\label{fig:academic_swelling_bark}
	\end{subfigure}
	\caption{Baker 算例下晶内气泡演化与肿胀行为验证}
	\label{fig:bark}
\end{figure}

\begin{figure}[!htb]
	\centering
	\begin{subfigure}[t]{0.45\textwidth}
		\centering
		\includegraphics[width=\linewidth,height=6.75cm,keepaspectratio]{academic_fgr_white.png}
		\caption{SCIANTIX 与 Miracle 裂变气体释放分数对比}
		\label{fig:academic_fgr_white}
	\end{subfigure}
	\hfill
	\begin{subfigure}[t]{0.45\textwidth}
		\centering
		\includegraphics[width=\linewidth,height=6.75cm,keepaspectratio]{academic_swelling_white.png}
		\caption{SCIANTIX 与 Miracle 晶界气体肿胀对比}
		\label{fig:academic_swelling_white}
	\end{subfigure}
	\caption{White 算例下晶界气泡演化与裂变气体释放行为验证}
	\label{fig:white}
\end{figure}

\begin{figure}[!htb]
	\centering
	\begin{subfigure}[t]{0.45\textwidth}
		\centering
		\includegraphics[width=\linewidth,height=6.75cm,keepaspectratio]{academic_fgr_hbs.png}
		\caption{高燃耗下 SCIANTIX 与 Miracle 裂变气体释放分数对比}
		\label{fig:academic_fgr_hbs}
	\end{subfigure}
	\hfill
	\begin{subfigure}[t]{0.45\textwidth}
		\centering
		\includegraphics[width=\linewidth,height=6.75cm,keepaspectratio]{academic_swelling_hbs.png}
		\caption{高燃耗下 SCIANTIX 与 Miracle 气体肿胀对比}
		\label{fig:academic_swelling_hbs}
	\end{subfigure}
	\caption{高燃耗结构重构与 HBS 孔隙演化行为验证}
	\label{fig:hbs}
\end{figure}

\begin{figure}[!htb]
	\centering
	\includegraphics[width=9.00cm,height=5.75cm,keepaspectratio]{academic_deviation.png}
	\caption{燃料化学计量偏差演化验证}
	\label{fig:academic_deviation}
\end{figure}

\xsubsection{间隙传热模型的数值测试与响应分析}{Numerical Testing and Response Analysis of the Gap Heat Transfer Model}
燃料芯块与包壳之间的间隙传热是燃料棒热力耦合分析中的关键环节，其换热能力受间隙宽度、气体组成、气体温度和腔体压力等因素共同影响。为验证本文实现的间隙传热模型及其与气体压力计算模块之间的耦合关系，本节构建二维接触传热算例，对不同间隙传热处理方式下的换热系数演化规律进行对比分析。

验证模型由两个尺寸均为 $1\times1$ 的矩形区域组成，两区域之间设置初始间隙距离为 $1$，用于表征燃料芯块与包壳之间的气体传热通道。计算中采用有限应变固体力学模型，并在接触界面施加基于 \texttt{GapHeatTransfer} 的间隙换热边界条件。间隙距离采用 \texttt{TOPTAN} 模型计算，气体热导率采用 \texttt{ADVANCED} 模型描述，以考虑气体状态变化对间隙传热能力的影响。

热边界条件设置如下：左侧边界施加随时间升高的温度载荷，温度由 $500~\mathrm{K}$ 逐渐升高至 $3000~\mathrm{K}$；右侧边界温度保持为 $400\,\mathrm{K}$，从而在接触界面两侧形成稳定的温度梯度。为突出气体间隙热阻对整体传热过程的控制作用，固体材料热导率设置为 $1.0\times10^8\,\mathrm{W\,m^{-1}\,K^{-1}}$，以尽量减小固体区域内部温降。力学边界条件中，对模型边界施加位移约束，以避免刚体位移对接触间隙计算产生干扰。

气体间隙初始充入 He，并通过 \texttt{CavityPressure} 模块计算腔体压力随温度变化的响应。计算过程中，间隙宽度、气体温度和间隙换热系数分别由辅助变量和后处理器提取，用于评价间隙传热模型对气体状态变化的响应能力。瞬态计算时间步长取 $0.1\,\mathrm{s}$，总计算时间为 $2\,\mathrm{s}$。

为分析压力反馈对间隙传热结果的影响，本文设置两组对比计算。第一组采用 \texttt{ADVANCED} 气体热导率模型，并耦合 \texttt{CavityPressure} 模块以考虑腔体压力随温度变化的反馈作用；第二组采用默认气隙导热处理方式，不显式考虑压力反馈对气体导热能力的影响。通过比较两组模型中间隙换热系数随温度升高的变化趋势，可以检验压力--温度--气体热导率耦合关系的实现正确性。

图 \ref{fig:gap_heat_transfer_verification} 给出了不同间隙传热模型下换热系数随温度变化的计算结果。随着边界温度升高，气体热导率逐渐增大，间隙换热系数整体呈上升趋势。当考虑腔体压力反馈时，气体温度升高引起压力和状态参数同步变化，使 \texttt{ADVANCED} 模型计算得到的气体导热能力进一步增强，因此对应的间隙换热系数高于默认模型结果。相比之下，默认模型未考虑压力和气体状态参数的动态反馈，其换热系数增长幅度相对较小。

\begin{figure}[H]
	\centering
	\includegraphics[width=9.00cm,height=5.0cm,keepaspectratio]{academic_verification_plot.png}
	\caption{不同间隙传热模型下换热系数随温度变化规律}
	\label{fig:gap_heat_transfer_verification}
\end{figure}

上述结果表明，本文实现的间隙传热模型能够正确反映气体温度、压力及间隙状态变化对界面换热能力的预期影响。两种模型所得到的变化趋势与对应模型的理论关系一致，说明程序中实现的间隙换热模块变量传递和耦合计算过程具有较好的物理合理性和数值稳定性，可为后续燃料棒尺度热--力--气体耦合计算提供数值基础。

\xsection{本章小结}{Summary}

本章针对燃料元件非线性力学、裂变气体行为与燃料–包壳间隙传热的数值耦合问题，通过材料单元测试、裂变气体基准算例及界面传热响应分析，评价了相关模型的数值实现及其物理响应特征。主要结论如下。

隐式本构积分能够较准确地再现燃料与包壳材料的时间相关非弹性变形行为。在恒温、恒载荷条件下，HT9、SS316和D9包壳蠕变模型的计算结果与相应理论参考解保持良好一致。UO₂燃料蠕变单元测试中，计算轴向位移与理论参考值的相对差异约为1.23\%；MOX燃料包含加载、应力松弛及二次加载阶段的算例中，第一阶段蠕变应变的相对差异约为1.05\%。上述结果表明，基于径向返回算法的局部本构积分能够再现相应材料模型规定的非弹性响应，并与全局力学平衡求解保持一致，为燃料与包壳蠕变、塑性及相关应力演化计算提供了数值验证依据。

裂变气体释放对扩散动力学及晶界饱和条件具有显著依赖，不同机理的时间响应需要相应的数值处理。ForMas模型的基准计算再现了晶界气体逐步积累、达到饱和阈值后释放及再次积累的周期性行为，累计释放量呈阶跃式增长。在相同热工况下，采用较大扩散系数的Turnbull模型会使晶界更早达到饱和条件，增加释放事件次数及累计释放量；相比之下，单次释放比例主要影响释放幅度和频率，对最终累计释放量的影响相对较弱。SCIANTIX耦合算例进一步表明，晶内气泡演化、晶界释放、高燃耗结构及热化学状态等主要响应能够与独立基准结果保持一致，但化学计量偏差的快速演化对时间步长较为敏感，适当减小时间步长能够改善数值结果。上述结果说明，裂变气体模型的计算表现不仅取决于扩散和释放机制，还受到内部状态变量更新及时间离散精度的影响。

间隙传热能力随气体温度和压力状态变化而改变，固定或简化气体状态处理可能低估相应工况下的换热增强。间隙传热算例表明，随着气体温度升高，气体热导率及间隙换热系数总体呈上升趋势。当考虑腔体压力反馈及气体状态动态更新时，ADVANCED模型预测的间隙换热系数高于未显式考虑相关反馈的默认模型，表明压力及气体状态变化会对界面换热产生可辨识的影响。结合裂变气体释放对气体组成和棒内压力的作用，可以看出，间隙传热不宜作为与运行历史无关的固定边界条件，而应随气体状态和燃料–包壳间隙状态动态更新。相关算例验证了所实现模型的预期响应规律，但其对整棒温度预测的综合影响仍需通过后续燃料棒级实验验证加以评价。

